\chapter*{Laboratorio matemático de pizarra: estimar y comprobar}
\addcontentsline{toc}{chapter}{Laboratorio matemático de pizarra}

\section*{Un experimento completo desde cero: cuatro símbolos en prueba}
\addcontentsline{toc}{section}{Un experimento completo desde cero}
Fijemos el alfabeto $V=\{a,b,c\}$, una cadena de entrenamiento
$\texttt{abacabbacb}$ y la cadena reservada $\texttt{abca}$.
Ambas son \textbf{ILUSTRATIVAS}: se construyeron para enseñar la cuenta, no provienen de
un corpus observado. No obstante, sus resultados numéricos están regenerados por
\texttt{codigo/s01.py} y contrastados en \texttt{verificaciones/}.

Construimos la matriz de transiciones $C_{ij}=\#\{t:x_t=i,x_{t+1}=j\}$:
\[
C=\begin{array}{c|rrr|r}&a&b&c&N_i\\\hline
 a&0&2&2&4\\ b&2&1&0&3\\ c&1&1&0&2\end{array}.
\]
Hay $10-1=9$ transiciones, no diez: el símbolo inicial no tiene predecesor.
La ecuación de normalización para $\alpha=1$ es
\[
 P_{ij}=\frac{C_{ij}+1}{N_i+3}
 =\begin{pmatrix}1/7&3/7&3/7\\1/2&1/3&1/6\\2/5&2/5&1/5\end{pmatrix}.
\]
El denominador de la primera fila es $7$, no $6$; el de la tercera es $5$.
La comprobación de suma por fila es elemental pero obligatoria:
$1/7+3/7+3/7=1$, $1/2+1/3+1/6=1$ y $2/5+2/5+1/5=1$.
La probabilidad \emph{condicional de las tres transiciones de prueba}, dado $X_1=a$, es
\[
 P(b\mid a)P(c\mid b)P(a\mid c)
 =\frac37\cdot\frac16\cdot\frac25=\frac1{35}.
\]
No es todavía la probabilidad de la cadena completa. Para ello hay que precisar
la distribución inicial $\pi$: si elegimos $\pi=(1/3,1/3,1/3)$,
\[
 q(\texttt{abca})=\pi(a)\,P(b\mid a)P(c\mid b)P(a\mid c)=\frac1{105}.
\]
Por tanto, la longitud ideal de código es $\log_2(105)\approx6.714246$ bits,
y, con cuatro eventos de predicción, la tasa es
$\log_2(105)/4\approx1.678561$ bits por carácter y la perplejidad
$105^{1/4}\approx3.201086$.
Si en lugar de cuatro eventos se evalúan solo las tres transiciones,
la pérdida media pasa a $\log_2(35)/3$; \emph{ambas cifras son válidas,
pero corresponden a experimentos condicionales diferentes}.

\begin{ejemplo}[Un contraste que detecta un modelo incorrecto]
El estimador no suavizado asigna $\widehat P(c\mid b)=0$.
Entonces $q(\texttt{abca})=0$ para cualquier $\pi(a)>0$ y la pérdida
es $+\infty$. Sustituir el cero únicamente al mirar la prueba sería
modificar el modelo después de observarla: hay que definir el suavizado y
seleccionar su intensidad en validación.
\end{ejemplo}

\begin{figure}[H]\centering
\begin{tikzpicture}[>=Stealth,font=\small, every node/.style={align=center}]
\node[draw=azul,fill=azulclaro,rounded corners,minimum width=3cm,minimum height=1cm] (t) at (0,0) {\textbf{Entrenamiento}\\10 símbolos};
\node[draw=azul,fill=azulclaro,rounded corners,minimum width=3cm,minimum height=1cm] (c) at (4,0) {\textbf{Conteos}\\9 transiciones};
\node[draw=verde,fill=verdeclaro,rounded corners,minimum width=3cm,minimum height=1cm] (p) at (8,0) {\textbf{Suavizado}\\filas de $P$};
\node[draw=ocre,fill=ocreclaro,rounded corners,minimum width=3cm,minimum height=1cm] (e) at (8,-2) {\textbf{Prueba}\\$\texttt{abca}$};
\node[draw=azul,fill=azulclaro,rounded corners,minimum width=3cm,minimum height=1cm] (r) at (2,-2) {\textbf{Evaluación}\\$\log_2 105$ bits};
\draw[->,thick] (t)--(c);\draw[->,thick] (c)--(p);\draw[->,thick] (p)--(e);\draw[->,thick] (e)--(r);
\end{tikzpicture}
\caption{La prueba solo recibe un modelo previamente fijado. El flujo contrario constituye fuga de información.}
\end{figure}

\section*{Regla de la cadena frente a memoria finita: una cuenta contrastable}
\addcontentsline{toc}{section}{Regla de la cadena y memoria finita}
Tomemos $X_1$ y $X_2$ Bernoulli$(1/2)$ independientes, y definamos
$X_3=X_1$ casi seguramente. Se trata de una distribución conjunta legítima,
aunque no sea markoviana de orden uno. Para $x=(1,0,1)$,
\[
 p(1,0,1)=p(X_1=1)\,p(X_2=0\mid X_1=1)\,
 p(X_3=1\mid X_1=1,X_2=0)=\frac12\frac12\cdot1=\frac14.
\]
Un modelo construido con los \emph{verdaderos} condicionales de orden uno
utilizaría $p(X_3=1\mid X_2=0)=1/2$ y asignaría $1/8$ a esa cadena.
Esta diferencia no es error algebraico: proviene de descartar información.

\begin{figure}[H]\centering
\begin{tikzpicture}[>=Stealth,font=\small,
 v/.style={draw=azul,fill=azulclaro,circle,minimum size=10mm}]
\node[v] (x1) at (0,0) {$X_1$};\node[v] (x2) at (4,0) {$X_2$};\node[v] (x3) at (8,0) {$X_3$};
\draw[->,thick] (x1)--(x2);\draw[->,thick] (x2)--(x3);
\draw[->,thick,dashed,color=vino] (x1) to[bend left=35] node[above] {memoria relevante} (x3);
\node at (4,-1.0) {Un bigrama descarta la dependencia discontinua $X_1\rightarrow X_3$.};
\end{tikzpicture}
\caption{Diagrama explicativo de una dependencia que no cabe en un único símbolo de memoria.}
\end{figure}

\section*{Suavizado como regularización: el papel del tamaño del vocabulario}
\addcontentsline{toc}{section}{Suavizado y tamaño del vocabulario}
Sea $K$ el tamaño del alfabeto y $\alpha>0$. El peso total de los
pseudoconteos es $K\alpha$, de modo que
\[
\frac{c_i+\alpha}{N+K\alpha}
=\underbrace{\frac{N}{N+K\alpha}}_{\lambda}\frac{c_i}{N}
 +\underbrace{\frac{K\alpha}{N+K\alpha}}_{1-\lambda}\frac1K.
\]
Esta identidad describe el suavizado como mezcla convexa entre el estimador
empírico y la distribución uniforme. Es válida para $N>0$ y $\alpha>0$.
Para $N=10$, $K=5$ y $\alpha=1$, $\lambda=2/3$; si se aumenta
$K$ a 100, manteniendo $N$ y $\alpha$, $\lambda=1/11$. La masa que
se introduce por suavizado depende fuertemente del vocabulario.

\begin{ejemplo}[Una decisión que exige validación]
En un contexto con conteos $(6,3,1,0,0)$, $K=5$ y $\alpha=1$,
se obtiene $q=(7,4,2,1,1)/15$. El símbolo ausente recibe $1/15$.
Con $\alpha=1/2$, el denominador es $12.5$ y recibe $0.5/12.5=1/25$.
Ninguna de las dos opciones es automáticamente superior: la pérdida en
validación sirve para decidir antes de examinar prueba.
\end{ejemplo}

\section*{Perplejidad: tres lecturas compatibles y una lectura incorrecta}
\addcontentsline{toc}{section}{Perplejidad y medias geométricas}
Fijadas las probabilidades $r_t=q(x_t\mid h_t)>0$ de $T$ eventos observados,
\[
\operatorname{PPL}
=\exp\left(-\frac1T\sum_{t=1}^T\ln r_t\right)
=\left(\prod_{t=1}^T\frac1{r_t}\right)^{1/T}.
\]
En consecuencia, la PPL es la \emph{media geométrica de las inversas}
de las probabilidades asignadas a los eventos. Si todas son iguales a
$1/K$, la PPL es $K$; para probabilidades variables, la PPL no es
el número real de símbolos candidatos ni la inversa de la exactitud.

\begin{ejemplo}[Producto, bits y exponencial: cálculo completo]
Para probabilidades observadas $(1/2,1/4,1/8,1/2)$,
\[
-\sum_{t=1}^4\log_2r_t=1+2+3+1=7\text{ bits},
\qquad\operatorname{PPL}=2^{7/4}\approx3.363586.
\]
Si la tercera probabilidad se reduce de $1/8$ a $1/16$ sin alterar las demás,
la suma aumenta a $8$ bits y la PPL a $2^2=4$. El cambio puede
explicarse multiplicativamente: se duplicó el código ideal de ese evento.
\end{ejemplo}

\section*{Comparar tasas, no solo nombres de métricas}
\addcontentsline{toc}{section}{Comparación justa entre tokenizaciones}
Dos modelos evalúan exactamente la misma cadena de $100$ caracteres,
sin pérdida de información al tokenizar. El modelo $A$ predice
$100$ caracteres con PPL $3.2$, mientras que $B$ predice $30$ subpalabras
con PPL $25$. Sus longitudes ideales son
\[
B_A=100\log_2(3.2)\approx167.8072\text{ bits},\quad
B_B=30\log_2(25)\approx139.3157\text{ bits}.
\]
Si el texto de referencia y el cálculo probabilístico sobre \emph{todo} el texto
son comparables, las tasas respectivas son $1.6781$ y $1.3932$ bits por
carácter. En este ejemplo el orden por PPL es inverso al orden por tasa.

\begin{nota}[Limitación esencial]
La misma unidad de longitud no convierte esta comparación en un efecto
\emph{causal} de la tokenización: también pueden variar datos de entrenamiento,
normalización, parametrización, contexto y presupuesto computacional. En particular,
no se identifica con BPC la calidad semántica de un modelo.
\end{nota}
Tampoco debe contarse el mismo dato dos veces: si una cadena Unicode se
normaliza antes del modelado, el número de caracteres del denominador debe
reflejar una convención fija y documentada. Para comparar código sin pérdida
se requiere poder reconstruir exactamente la cadena original.

\section*{De muestras a poblaciones: qué se demuestra realmente}
\addcontentsline{toc}{section}{Diferencia entre riesgo empírico y poblacional}
La pérdida observada sobre $T$ eventos es
$\widehat R_T(q)=T^{-1}\sum_{t=1}^T-\ln q(x_t\mid h_t)$.
En un problema categórico \emph{i.i.d.} con observaciones de distribución
$p$ y soporte compatible, la esperanza de $-\ln q(X)$ es $H(p,q)$;
la ley de los grandes números justifica convergencia para $q$ fijo
si la variable es integrable. En una secuencia dependiente la convergencia
requiere hipótesis apropiadas (por ejemplo, ergodicidad e integrabilidad),
y en modelos elegidos usando los mismos datos aparece un problema adicional
de sobreajuste.

Por otro lado, para una muestra fija, minimizar $\widehat R_T(q_\theta)$
es exactamente maximizar la verosimilitud correspondiente, porque
\[
\arg\min_\theta\widehat R_T(q_\theta)
=\arg\min_\theta\left(-\frac1T\ln q_\theta(x_{1:T})\right)
=\arg\max_\theta q_\theta(x_{1:T}).
\]
Esto no prueba que $q_\theta=p$, ni establece por sí solo generalización.
El paso de lo observado a una afirmación científica exige particiones,
incertidumbre y una pregunta experimental bien definida.
